% =====================================================================
%  control_cell.tex  --  top view of the two control-cell types, coloured as in
%  the MCNP render (yellow water, green aluminium, red graphite).  Left: a fully
%  graphite-filled control cell (no hole).  Right: a holed control cell, where a
%  central water hole (bore d = 2.9) is lined by an aluminium wall (w = 1.5 mm),
%  giving an outer diameter of 3.2; only this hole is dimensioned (all cm).  Both
%  diameters are placed on the right of the holed cell (outer high, inner low) so
%  they clear the filled cell on the left.  The cells use the same coordinate unit
%  (0.524 cm) as fuel_cell_types.tex, so the squares are the same size.
%  \input from chapters/2-geometry.tex.
% =====================================================================
\begin{figure}[H]
  \centering
  \begin{tikzpicture}[
      x=0.524cm,y=0.524cm,
      font=\small,
      lead/.style={-{Stealth[length=1.4mm]}, thin},
    ]
    % phantom points: balance the bounding box about the pair midpoint (x=4.4) so
    % the two cells stay centred despite the holed cell's right-side dimensions.
    \path (-5.5,0) (14.3,0);
    % ===== left: fully filled graphite control cell (no hole) =====
    \begin{scope}[shift={(0,0)}]
      \fill[yellow]         (-3.6,-3.6)   rectangle (3.6,3.6);    % water (cell)
      \fill[green!75!black] (-3.4,-3.4)   rectangle (3.4,3.4);    % Al box wall
      \fill[red]            (-3.25,-3.25) rectangle (3.25,3.25);  % graphite (cavity)
      \draw[thick] (-3.6,-3.6)   rectangle (3.6,3.6);     % cell pitch
      \draw        (-3.4,-3.4)   rectangle (3.4,3.4);     % Al box outer wall
      \draw        (-3.25,-3.25) rectangle (3.25,3.25);   % Al box cavity
    \end{scope}
    % ===== right: holed control cell (dimensioned) =====
    \begin{scope}[shift={(8.8,0)}]
      \fill[yellow]         (-3.6,-3.6)   rectangle (3.6,3.6);    % water (cell)
      \fill[green!75!black] (-3.4,-3.4)   rectangle (3.4,3.4);    % Al box wall
      \fill[red]            (-3.25,-3.25) rectangle (3.25,3.25);  % graphite (cavity)
      \fill[green!75!black] (0,0) circle (1.6);                   % Al wall (D = 3.2)
      \fill[yellow]         (0,0) circle (1.45);                  % water hole (D = 2.9)
      \draw[thick] (-3.6,-3.6)   rectangle (3.6,3.6);     % cell pitch
      \draw        (-3.4,-3.4)   rectangle (3.4,3.4);     % Al box outer wall
      \draw        (-3.25,-3.25) rectangle (3.25,3.25);   % Al box cavity
      \draw        (0,0) circle (1.6);
      \draw        (0,0) circle (1.45);
      % ---- diameter dimensions on the right: outer (D 3.2) high, inner (D 2.9) low ----
      \draw[lead] (3.9,1.8)  -- (1.453,0.671);            % D = 3.2 (Al wall, outer)
      \node[anchor=west,fill=white,inner sep=1.5pt] at (3.95,1.8) {\O 3.2};
      \draw[lead] (3.9,-1.8) -- (1.317,-0.608);           % D = 2.9 (water bore, inner)
      \node[anchor=west,fill=white,inner sep=1.5pt] at (3.95,-1.8) {\O 2.9};
    \end{scope}
  \end{tikzpicture}
  \caption{Top view of the two control-cell types.}
  \label{fig:controlcell}
\end{figure}
